\section{Geometry, charge conventions and scope}
\label{sec:probe-setup}

Magnetic objects can test aspects of a Calabi--Yau transition
that are not fixed by its massless electric spectrum. Such a
test requires both a charge calculation and a stability
condition. This note studies the first requirement and parts
of the second for the compact transition of
\cite{Wuebben:2026quantum}. We compute specified numerical
Donaldson--Thomas contributions, their polarization
dependence, a two-constituent magnetic wall near large volume,
and exact phase constraints near the quantum transition.
The results do not determine the full physical BPS index at
the strongly coupled endpoint.

Let $X=\wideX_9$ be the smooth resolution of a compact
Calabi--Yau threefold with two conifold points and one cone
over $E=\operatorname{Bl}_2\mathbb P^2$. Its Hodge numbers
are $(6,122)$. The nodal curves $C_1,C_2$ and the two
exceptional curves $e_1,e_2$ of $E$ obey
\begin{equation}
 R=\iota_*(e_1-e_2)=C_1+C_2.
 \label{eq:probe-root}
\end{equation}
The compact gaugings allow these local sectors to Higgs only
together. The smooth side has Hodge numbers $(3,123)$ and a
rigid plane. In the resolved chamber the corresponding
surface is $S_-=D_8\simeq\mathbb F_1$; on the smooth side
it is $S_+\simeq\mathbb P^2$, with class $\nu$.
Their classical restrictions and string anomalies are computed
in \cite[Section~3]{Wuebben:2026quantum}.

The geometric transition exists both by an explicit construction and by
the exact mixed smoothing criterion
\cite[Theorems~4.7 and~11.8]{Wuebben:2026smoothing}. The relation
\eqref{eq:probe-root} has nonzero coefficients at all three germs, and
the selected deformation base of the singular threefold is smooth.
That theorem deforms the contracted threefold, not a sheaf on its
resolution. To avoid conflating the two problems, we prove below a
separate persistence statement for the isolated sheaves under smooth
polarized deformations of $X$; transport across the topology-changing
transition remains outside that assertion.

We use the primitive divisor basis
\begin{equation}
 \mathcal D=(D_2,D_5,D_6,D_7,D_8,E).
 \label{eq:x9-divisor-basis}
\end{equation}
The integral nef basis $(H_1,\ldots,H_6)$ has columns
\begin{equation}
 \begin{pmatrix}
 0&0&1&1&2&2\\1&1&2&4&5&5\\0&1&2&4&4&5\\
 0&-1&-1&-3&-3&-3\\0&-1&-1&-2&-2&-2\\0&0&1&2&2&2
 \end{pmatrix}
 \label{eq:x9-mirror-nef}
\end{equation}
in $\mathcal D$. In particular
$H_1=D_5$, $H_2=D_4+D_5$, $H_3=D_1$, $H_4=D_0$.
Write $\beta_a$ for the dual curve basis. Then
\begin{equation}
 C_2=\beta_1,\quad C_1=\beta_5,\quad
 \iota_*e_2=\beta_6,\quad
 \iota_*e_1=\beta_1+\beta_5+\beta_6.
\end{equation}
The ample class $J_*=
\sum_aH_a=(6,18,16,-11,-8,7)$ fixes the initial polarization.
The contraction face is $J=aD_0+bD_1$, $a,b\geq0$.
Appendix~\ref{app:probe-data} prints the intersection data
needed for the sheaf stability tests.

For a pure surface sheaf $\mathcal F$, charges mean
\begin{equation}
 \Gamma(\mathcal F)=\operatorname{ch}(\mathcal F)
          \sqrt{\operatorname{Td}(X)}=(0,P,q,q_0),\qquad
 Z_{\rm LV}=-\int_Xe^{-t}\Gamma,\quad t=B+iJ.
\end{equation}
The D0 convention makes a point sheaf have period $-1$.
The Dirac pairing with an electric charge
$(0,0,\gamma,\gamma_0)$ is $P\gamma$.
The normalized exact central charge for $P\neq0$ differs
from its polynomial value by a magnetic period correction
$\Psi_P$. No additional curvature term is added after the
Mukai charge has been fixed.

The four-dimensional particles here come from wrapping a
five-dimensional magnetic string on an auxiliary spacetime
circle. This is distinct from the elliptic circle in the
six-dimensional interpretation of $X$. The physical mass
contains the usual common normalization; phase comparisons
use the holomorphic central charges above.

The main unconditional sheaf result is the existence of two
isolated stable components carrying relative charges
$(C_2,-1)$ and $(R,-1)$ with respect to a moving D4 seed.
Each contributes $+1$ to the numerical DT invariant.
Both remain Gieseker stable along ample polarizations approaching
the contraction. Their explicitly identified magnetic split
has a transverse alignment near large volume, but the exact
endpoint periods lie strictly away from that alignment.
Light-electric emission walls are also excluded on a final
specified segment. These exclusions are incomplete: they do
not prove that the constituent BPS states survive the whole
continuation. Moreover their magnetic charges are confined at
a generic Higgs vacuum. Conditional counts involving the full
moving pencils are kept in Appendix~\ref{app:pencil} so that
their additional genericity assumption is separate from the
isolated-component results.
